\Def Пифагорова тройка: $(a, b, c): a, b, c \in \N; a^2 + b^2 = c^2$

\Note Можно рассматривать только взаимно простые числа: если любые два числа из тройки имеют общий множитель, то и третье число делится на него.

\Th 
Если $a^2 + b^2 = c^2$ для попарно взаимно простых чисел $a, b$, то с точностью до перестановки a и b и их знака найдутся такие целые числа $m, n$, что $a = m^2 - n^2, b = 2mn, c = m^2 + n^2$

\Proof

1. Перепишем исходное равенство следующим образом: $a^2 + b^2 = (a+bi)(a-bi) = c^2$. 

2. $a+bi$ и $a - bi$ взаимно просты: $(a+bi, a-bi) = (a+bi, 2a)$. Заметим, что $(a + bi, 2) = 1$, т.к. иначе $c^2$ кратно 2, а значит, и 4, а значит, и a, и b, кратны 2 (т.к. иначе $c^2$ либо нечётное, либо кратно 2, но не кратно 4), а это противоречит взаимной простоте $a, b, c$. Тогда $(a+bi, a-bi) = (a+bi, 2a) = (a+bi, a) = (bi, a) = (a, b) = 1$ в силу взаимной простоты.

3. Произведение $a+bi$ и $a - bi$ - квадрат какого-то числа, получается, всякое простое число, которое встречается в разложении $c^2$ (а оно встречается чётное число раз), встречается и в разложении либо $a + bi$, либо $a-bi$ то же число раз. Таким образом, $a + bi = u(m + ni)^2$, где $u$ - обратимый элемент. Сделав замену при необходимости, можно избавиться от u; Таким образом, $a + bi = (m + ni)^2$.

4. Тогда $a = m^2 - n^2, b = 2mn$; $a - bi = \overline{a + bi} = \overline{(m + ni)^2} = (\overline{m + ni})^2 = (m - ni)^2$; тогда $c = (m + ni)(m - ni) = m^2 + n^2$ с точностью до знака.
\EndProof